% 第4章: データの流れを整理する
\section{データの流れを整理する}

%% ── Q1: 直接アクセスの問題 ──────────────────
\needspace{5\baselineskip}
\begin{qabox}{モジュール同士が直接話すとなぜダメ？}

「温度が30度を超えたらLEDを点灯」したいとき、こんなコードを書きたくなるかもしれません。

\begin{badexample}{モジュール同士が直接参照}
\begin{lstlisting}
class TempSensorModule {
private:
    LedModule* _led;  // LEDモジュールを直接持っている！

public:
    void update() {
        float temp = readSensor();
        if (temp > 30.0) {
            _led->turnOn();  // 直接操作
        }
    }
};
\end{lstlisting}
\end{badexample}

一見うまく動きそうですが、以下の問題が起きます。

\begin{itemize}
  \item 温度センサーモジュールがLEDモジュールに依存 → LEDなしでは使えない
  \item 「温度でモーターも回したい」→ MotorModule*も追加？ → 依存がどんどん増える
  \item モジュール間の依存が複雑に絡み合い、変更が困難に（スパゲッティ化）
\end{itemize}

\begin{center}
\begin{tikzpicture}[
  module/.style={rectangle, draw=warnred!70, fill=warnred!5,
    minimum width=28mm, minimum height=10mm, font=\small\bfseries},
  arr/.style={-{Stealth[length=3mm]}, thick, warnred!70},
]
  \node[module] (temp) {TempSensor};
  \node[module, right=25mm of temp] (led) {LED};
  \node[module, below=12mm of led] (motor) {Motor};
  \node[module, left=25mm of motor] (disp) {Display};

  \draw[arr] (temp) -- (led);
  \draw[arr] (temp) -- (motor);
  \draw[arr] (led) -- (disp);
  \draw[arr] (motor) -- (disp);
  \draw[arr] (disp) to[bend right=20] (temp);
\end{tikzpicture}
\end{center}

矢印（依存関係）が蜘蛛の巣のように広がっています。これでは1つのモジュールを変更すると、他のモジュールにも影響が波及します。

\end{qabox}

%% ── Q2: SystemDataとは ──────────────────────
\needspace{5\baselineskip}
\begin{qabox}{SystemDataって何？}

\texttt{SystemData}は、すべてのモジュールのデータを1か所にまとめた「伝言板」です。

\begin{goodexample}{SystemData — モジュールの伝言板}
\begin{lstlisting}
#include "LedModule.h"
#include "TempSensorModule.h"

struct SystemData {
    LedData led;             // LEDの状態
    TempSensorData temp;     // 温度センサーの値
};
\end{lstlisting}
\end{goodexample}

モジュール同士は直接会話しません。代わりにSystemDataを介してデータをやり取りします。

\begin{center}
\begin{tikzpicture}[
  module/.style={rectangle, draw=accentblue, fill=accentblue!8,
    minimum width=28mm, minimum height=10mm, font=\small\bfseries},
  sysdata/.style={rectangle, draw=okgreen!70, fill=okgreen!8,
    minimum width=35mm, minimum height=14mm, font=\bfseries},
  arr/.style={-{Stealth[length=3mm]}, thick, accentblue},
]
  \node[sysdata] (sd) {SystemData};

  \node[module, above left=14mm and 15mm of sd] (temp) {TempSensor};
  \node[module, above right=14mm and 15mm of sd] (camera) {Camera};
  \node[module, below left=14mm and 15mm of sd] (led) {LED};
  \node[module, below right=14mm and 15mm of sd] (tft) {TFT};

  \draw[arr] (temp) -- node[left, font=\scriptsize]{書き込み} (sd);
  \draw[arr] (camera) -- node[right, font=\scriptsize]{書き込み} (sd);
  \draw[arr] (sd) -- node[left, font=\scriptsize]{読み取り} (led);
  \draw[arr] (sd) -- node[right, font=\scriptsize]{読み取り} (tft);
\end{tikzpicture}
\end{center}

\begin{itemize}
  \item 入力モジュール（TempSensor等）はSystemDataにデータを書き込む
  \item 出力モジュール（LED等）はSystemDataからデータを読み取る
  \item モジュール同士は直接参照しない → 依存関係がシンプル
\end{itemize}

温度センサーはLEDの存在を知りません。ただSystemDataの \texttt{temp} に値を書くだけです。LEDをモーターに差し替えても、温度センサーのコードは1行も変わりません。

\end{qabox}

%% ── Q3: 前方宣言 ──────────────────────────
\needspace{5\baselineskip}
\begin{qabox}{前方宣言って何？}

IModule.hの中にこんな1行がありました。

\begin{lstlisting}
struct SystemData;  // 前方宣言
\end{lstlisting}

これは「SystemDataという型がどこかで定義されてますよ」とコンパイラに教える宣言です。

\subsection{なぜ必要なの？}

IModule.hは \texttt{update(SystemData\& data)} を宣言していますが、SystemDataの中身を知る必要はありません（ポインタや参照で受け取るだけなので）。

もしIModule.hが \texttt{SystemData.h} をインクルードすると、こんな問題が起きます。

\begin{badexample}{循環依存}
\begin{lstlisting}
// IModule.h が SystemData.h をインクルード
// SystemData.h が LedModule.h をインクルード
// LedModule.h が IModule.h をインクルード
// → IModule.h が SystemData.h をインクルード → 無限ループ！
\end{lstlisting}
\end{badexample}

前方宣言を使えば、この循環を断ち切れます。

\begin{goodexample}{前方宣言で循環を回避}
\begin{lstlisting}
// IModule.h — SystemData.hをインクルードしない
struct SystemData;  // 「こういう型がある」とだけ宣言

class IModule {
public:
    virtual void update(SystemData& data) = 0;
    // 参照で受け取るだけなので、中身を知らなくてOK
};
\end{lstlisting}
\end{goodexample}

\begin{notebox}[ヘッダーとcppの使い分け]
\begin{itemize}
  \item \textbf{.hファイル} — \texttt{struct SystemData;} の前方宣言のみ
  \item \textbf{.cppファイル} — \texttt{\#include "SystemData.h"} でフル定義を使う
\end{itemize}
.cppでは実際にSystemDataのメンバにアクセスするので、フル定義が必要です。このルールを守れば循環依存は起きません。
\end{notebox}

\end{qabox}
